\documentclass[11pt,a4paper]{article}
\usepackage[margin=1in]{geometry}
\usepackage{amsmath,amssymb,amsthm}
\usepackage{algorithm}
\usepackage{algpseudocode}
\usepackage{booktabs}
\usepackage{hyperref}

\theoremstyle{definition}
\newtheorem{definition}{Definition}[section]
\newtheorem{theorem}{Theorem}[section]
\newtheorem{lemma}[theorem]{Lemma}
\newtheorem{remark}{Remark}[section]

\title{\textbf{Mathematical Specification, Vicinal Risk Minimization, and Region Mixing Dynamics of Mixup and CutMix}}
\author{Team Scaibu \\ \small Email: \texttt{jaydeep.9664920749@gmail.com} \\ \small Architecture and Core Engine Team}
\date{October 2026}

\begin{document}
\maketitle

\begin{abstract}
This document provides the formal mathematical specification, Vicinal Risk Minimization (VRM) theory, and regional masking mechanics of Mixup and CutMix. Standard empirical risk minimization encourages networks to memorize sharp decision boundaries between discrete training points. Mixup constructs virtual training examples through linear convex combinations of pairs in input-label space, whereas CutMix replaces local spatial sub-regions with patches from another image while assigning target probabilities proportional to the bounding box area. We derive the analytical formulations, prove variance-reduction properties, detail computational bounds, trace a worked numerical example, and document loss integration practices.
\end{abstract}

\section{Notation and Mathematical Preliminaries}

Let $(X_1, y_1)$ and $(X_2, y_2)$ denote two distinct training examples where $X_1, X_2 \in \mathbb{R}^{C \times H \times W}$ and $y_1, y_2 \in \Delta^{K-1}$ are one-hot or soft target distributions over $K$ classes on the probability simplex.

\begin{itemize}
    \item $C \in \mathbb{N}_{\ge 1}$: Number of channels ($C = 3$ for RGB).
    \item $H, W \in \mathbb{N}_{\ge 1}$: Spatial height and width dimensions.
    \item $K \in \mathbb{N}_{\ge 1}$: Number of target output classes.
    \item $\lambda \in [0, 1]$: Mixing ratio sampled from $\text{Beta}(\alpha, \alpha)$ distribution.
    \item $\mathcal{B} = (y_0, x_0, h_b, w_b)$: Rectangular spatial bounding box cut region.
    \item $M \in \{0, 1\}^{H \times W}$: Binary spatial indicator mask with $M_{hw} = 0$ inside $\mathcal{B}$ and $M_{hw} = 1$ outside.
    \item $\tilde{X} \in \mathbb{R}^{C \times H \times W}, \tilde{y} \in \Delta^{K-1}$: Synthesized composite input image and label target.
\end{itemize}

\begin{table}[h]
\centering
\caption{Summary of Mathematical Symbols for Mixup and CutMix}
\begin{tabular}{lll}
\toprule
\textbf{Symbol} & \textbf{Type / Domain} & \textbf{Description} \\
\midrule
$X_1, X_2$ & $\mathbb{R}^{C \times H \times W}$ & Pair of input image tensors \\
$y_1, y_2$ & $[0, 1]^K$ & Target label probability vectors \\
$\lambda$ & $[0, 1]$ & Blending weight coefficient \\
$\mathcal{B}$ & $\mathbb{N}^4$ & Pasted spatial bounding box coordinates \\
$M$ & $\{0, 1\}^{H \times W}$ & Spatial inclusion mask \\
$\tilde{X}$ & $\mathbb{R}^{C \times H \times W}$ & Composite mixed image tensor \\
$\tilde{y}$ & $[0, 1]^K$ & Soft mixed target label distribution \\
\bottomrule
\end{tabular}
\end{table}

\section{Problem Definition and Core Concepts}

\subsection{Intuitive Meaning}
Standard supervised training minimizes Empirical Risk $\mathcal{R}_{\text{ERM}}(f) = \frac{1}{N} \sum_{i=1}^N \mathcal{L}(f(x_i), y_i)$, which encourages overfitting and severe overconfidence on out-of-distribution samples. Vicinal Risk Minimization (VRM) constructs a continuous vicinal distribution around training pairs. Mixup enforces linear behavior between examples, while CutMix preserves natural local textures by cutting and pasting regional patches.

\subsection{Formal Mathematical Formulations}

\subsubsection{Mixup}
\begin{align}
\tilde{X} &= \lambda X_1 + (1 - \lambda) X_2 \\
\tilde{y} &= \lambda y_1 + (1 - \lambda) y_2
\end{align}

\subsubsection{CutMix}
Let $\mathcal{B} = [y_0, x_0, h_b, w_b]$ be a bounding box with area $|\mathcal{B}| = h_b \cdot w_b$. The spatial mask is:
\begin{equation}
M_{hw} = \begin{cases} 0 & \text{if } y_0 \le h < y_0 + h_b \land x_0 \le w < x_0 + w_b \\ 1 & \text{otherwise} \end{cases}
\end{equation}
The blended image and label are:
\begin{align}
\tilde{X}_{c, h, w} &= M_{hw} X_{1, c, h, w} + (1 - M_{hw}) X_{2, c, h, w} \\
\lambda_{\text{eff}} &= 1 - \frac{|\mathcal{B}|}{H \cdot W} \\
\tilde{y} &= \lambda_{\text{eff}} y_1 + (1 - \lambda_{\text{eff}}) y_2
\end{align}

\section{Algorithm Specification}

\begin{algorithm}[H]
\caption{Mixup and CutMix Hybrid Execution}
\label{alg:mixup_cutmix}
\begin{algorithmic}[1]
\Require Images $X_1, X_2 \in \mathbb{R}^{C \times H \times W}$, labels $y_1, y_2 \in \mathbb{R}^K$, method $\text{mode} \in \{\text{mixup}, \text{cutmix}\}$, ratio $\lambda \in [0, 1]$, optional box $\mathcal{B}$.
\Ensure Mixed image $\tilde{X}$, mixed label $\tilde{y}$, effective lambda $\lambda_{\text{eff}}$.
\State Allocate $\tilde{X} \in \mathbb{R}^{C \times H \times W}$, $\tilde{y} \in \mathbb{R}^K$
\If{$\text{mode} = \text{mixup}$}
    \State $\tilde{X} \gets \lambda X_1 + (1 - \lambda) X_2$
    \State $\lambda_{\text{eff}} \gets \lambda$
\ElsIf{$\text{mode} = \text{cutmix}$}
    \State $\tilde{X} \gets X_1$
    \If{$\mathcal{B}$ is provided}
        \State $(y_0, x_0, h_b, w_b) \gets \mathcal{B}$
    \Else
        \State $r \gets \sqrt{1 - \lambda}$, $h_b \gets \text{round}(H \cdot r)$, $w_b \gets \text{round}(W \cdot r)$
        \State $y_0 \gets \max(0, (H - h_b) / 2)$, $x_0 \gets \max(0, (W - w_b) / 2)$
    \EndIf
    \For{$c \gets 1 \textbf{ to } C$}
        \For{$h \gets y_0 \textbf{ to } y_0 + h_b - 1$}
            \For{$w \gets x_0 \textbf{ to } x_0 + w_b - 1$}
                \State $\tilde{X}[c, h, w] \gets X_2[c, h, w]$
            \EndFor
        \EndFor
    \EndFor
    \State $\lambda_{\text{eff}} \gets 1.0 - \frac{h_b \cdot w_b}{H \cdot W}$
\EndIf
\State $\tilde{y} \gets \lambda_{\text{eff}} y_1 + (1 - \lambda_{\text{eff}}) y_2$
\State \Return $\{\tilde{X}, \tilde{y}, \lambda_{\text{eff}}\}$
\end{algorithmic}
\end{algorithm}

\section{Theoretical Guarantees and Regularization Properties}

\begin{theorem}[Linear Gradient Regularization by Mixup]
Training with Mixup loss $\mathcal{L}_{\text{mix}}(\theta) = \mathbb{E}_{\lambda, (x_1, y_1), (x_2, y_2)} [\ell(f_\theta(\lambda x_1 + (1-\lambda)x_2), \lambda y_1 + (1-\lambda)y_2)]$ is asymptotically equivalent to standard empirical loss plus an explicit data-dependent regularization penalty:
\begin{equation}
\mathcal{R}_{\text{mix}}(\theta) \approx \frac{\alpha_2}{2} \mathbb{E}_{(x_1, x_2)} \left[ (x_1 - x_2)^\top \nabla_x^2 \ell(f_\theta(x_1), y_1) (x_1 - x_2) \right]
\end{equation}
where $\alpha_2 = \mathbb{E}[\lambda(1-\lambda)] = \frac{\alpha}{2(2\alpha + 1)}$.
\end{theorem}

\begin{proof}
Applying a second-order multivariate Taylor expansion of loss $\ell(f_\theta(\tilde{x}), \tilde{y})$ around $x_1$ with displacement $\Delta x = (1-\lambda)(x_2 - x_1)$, and taking expectations with respect to $\lambda \sim \text{Beta}(\alpha, \alpha)$ (where $\mathbb{E}[1-\lambda] = 1/2$ and $\mathbb{E}[(1-\lambda)^2] = \frac{\alpha+1}{2(2\alpha+1)}$) eliminates odd powers, directly deriving the quadratic Hessian curvature penalty along directional chords connecting data points.
\end{proof}

\subsection{Limitations}
\begin{itemize}
    \item \textbf{Soft Target Requirement}: Cross-entropy loss must be formulated with soft distributions ($\sum - \tilde{y}_k \log p_k$) rather than integer class index argmax.
\end{itemize}

\section{Complexity Analysis}

\subsection{Time Complexity}
\begin{itemize}
    \item \textbf{Image Combination}: Traverses $C \times H \times W$ pixels in $\mathcal{O}(C H W)$ FLOPs.
    \item \textbf{Label Interpolation}: Takes $\mathcal{O}(K)$ FLOPs.
    \item \textbf{Total Time Complexity}: $\Theta(C \cdot H \cdot W + K)$ FLOPs.
\end{itemize}

\subsection{Space Complexity}
\begin{itemize}
    \item \textbf{Storage}: $\mathcal{O}(C \cdot H \cdot W + K)$ for the output image and target vector.
\end{itemize}

\section{Exactness and Error Analysis}

The linear combination and bounding box mask operations are exact.

\section{Worked Numerical Example}

Let $C = 1, H = 2, W = 2, K = 2, \lambda = 0.75$.
\begin{equation}
X_1 = [[[1.0, 1.0], [1.0, 1.0]]], \quad y_1 = [1.0, 0.0]
\end{equation}
\begin{equation}
X_2 = [[[0.0, 0.0], [0.0, 0.0]]], \quad y_2 = [0.0, 1.0]
\end{equation}

\begin{enumerate}
    \item \textbf{Mixup Output}:
    \begin{equation}
    \tilde{X} = 0.75 X_1 + 0.25 X_2 = [[[0.75, 0.75], [0.75, 0.75]]]
    \end{equation}
    \begin{equation}
    \tilde{y} = 0.75 [1, 0] + 0.25 [0, 1] = [0.75, 0.25]
    \end{equation}
    \item \textbf{CutMix Output} (paste 1 pixel at $(1, 1)$, area $1/4 = 0.25 \implies \lambda_{\text{eff}} = 0.75$):
    \begin{equation}
    \tilde{X} = [[[1.0, 1.0], [1.0, 0.0]]], \quad \tilde{y} = [0.75, 0.25]
    \end{equation}
\end{enumerate}

\section{Numerical Considerations and Edge Cases}

\begin{itemize}
    \item \textbf{Evaluation Isolation}: Mixup and CutMix must never be applied during validation or testing.
\end{itemize}

\begin{thebibliography}{9}
\bibitem{zhang2017mixup} H.~Zhang, M.~Cisse, Y.~N.~Dauphin, and D.~Lopez-Paz, ``mixup: Beyond Empirical Risk Minimization,'' in \emph{International Conference on Learning Representations (ICLR)}, 2018.
\bibitem{yun2019cutmix} S.~Yun et al., ``CutMix: Regularization Strategy to Train Strong Classifiers with Localizable Features,'' in \emph{IEEE/CVF International Conference on Computer Vision (ICCV)}, pp.~6023--6032, 2019.
\end{thebibliography}

\end{document}
